\documentclass[11pt]{article}
\usepackage[a4paper,margin=25mm]{geometry}
\usepackage{amsmath,amssymb,amsthm,booktabs,xurl,hyperref}
\hypersetup{hidelinks}
\setlength{\emergencystretch}{2em}
\newtheorem{theorem}{Theorem}
\newtheorem{lemma}[theorem]{Lemma}
\newtheorem{proposition}[theorem]{Proposition}
\newcommand{\Q}{\mathbb Q}
\newcommand{\Z}{\mathbb Z}
\title{A common arithmetic progression of contact arguments\\
and shared coefficients of their centered powers}
\author{\(\mathbb Z\mathbb Q\) Diophantine Project}
\date{23 November 2026}
\begin{document}
\maketitle
\begin{abstract}
The simultaneous general and specified integer certificates allow the
contact lists of different inputs to be translated after their point
differences are fixed. One sufficiently large spacing separates all
these lists by an arithmetic progression. We make that spacing a single
quantified coordinate and share the coefficients of the first contact's
centered square and fourth power. For twelve input blocks, seven common
coordinates replace thirty-three private coordinates. The complete
corrected quartic systems fall from $1074$ to $1048$ variables and from
$1254$ to $1228$ variables, with $908$ and $1064$ quadratic residuals.
Every omitted residual has an exact polynomial reconstruction identity.
Composition with the signed-binomial automatic positivity theorem
further gives $1043$ variables and $907$ rows, and $1223$ variables
and $1063$ rows, respectively.
The arithmetic assumptions of the existing interfaces are retained;
the corrected source uses the classical Jones coding theorem, without
the historical unresolved expanded-source premise.
\end{abstract}

\section{Scope and the order of the arithmetic choices}

We start from the corrected ceiling-power tables of
\cite{corrected,ceiling}. Their twelve integer-certificate inputs, in
the order used below, are
\[
 b,g,t,m,w,s,h,c,\gamma,o,j,i.
\]
The last input uses its existing reciprocal chart. The symbols in this
list belong to the outer source; the contact gauge $\rho$ used below is
a different common rational coordinate. The general interface has its
existing fixed defining coefficient $C=149^H$; the specified interface
has its previously fixed numerical coefficients.

All arithmetic assumptions of the global-gauge theorem \cite{global}
are retained: the full-domain reflected-contact height and detector
estimates, scalar and native-pole realization, simultaneous scalar and
point preparation, and the compatible finite-family prime supplier with
its all-sign order, finite avoidance, and approximation conclusions.
No $abc$ assumption or additional prime-supplier assertion is introduced.
The common-numerator theorem \cite{common} supplies the same compatible
index congruences as before.

In the simultaneous completeness construction, fix the integer inputs,
their root points, the common reference neighborhoods and point roles,
and their aggregate shifts. Each block then has a fixed finite set of
contact positions. In the general interface its two lists are internally
disjoint because their displacement is greater than seven while the
contact diameter is four. In the specified interface the retained two
successive shift bounds are seven and nine while the contact diameter
is $25/11$; these facts are exactly the separation facts used in
\cite{global}. Write $\mathcal A_i\subset\Q$ for the full internally
separated contact list of block $i$, with $0\le i<k$.

\begin{lemma}[One spacing separates a finite family]\label{spacing}
For any fixed finite sets $\mathcal A_0,\ldots,\mathcal A_{k-1}$,
there is a positive integer $D$ such that the sets
\[
                         \mathcal A_i-iD
\]
are pairwise disjoint. The block contact parameters may therefore be
chosen as $\lambda_i=\lambda+iD$ in the simultaneous completeness
construction.
\end{lemma}
\begin{proof}
Choose rational $L\le U$ containing the union of all the sets, and
take an integer $D>U-L$. The intervals
$[L-iD,U-iD]$ are pairwise disjoint. The roots, as functions of the
common slope $\lambda$, of the contact factors belonging to the
argument $\lambda+iD$ are precisely $\mathcal A_i-iD$.
\end{proof}

The spacing is chosen \emph{after} the aggregate shifts, and before
the contact primes. It need not have a uniform numerical upper bound
over all integer inputs. The equations below quantify its scaled or
unscaled value, so their coefficients do not depend on this choice.

Apply the existing finite-family supplier once to the disjoint union
in Lemma~\ref{spacing}. Enlarge finite avoidance by the supports of
these rational positions, their nonzero differences, the already fixed
point and scalar denominators, and all previous preparation moduli.
The common-index congruences and synchronized formal depths of
\cite{global,common} apply without change. At a selected contact only
its own weight has negative odd valuation; the other weights are
integral by the same finite avoidance. The global contact sum therefore
has its former valuation $-1$. At foreign contacts the old rational
square parity argument still applies.

At every previously prepared finite place and at real infinity,
$iD$ and all the aggregate displacements are now fixed. Thus the
existing reflected-contact expansion
\[
 q(\lambda+iD+u)
       =\kappa\lambda^{-2}+O(\lambda^{-3}),\qquad \kappa>0,
\]
holds simultaneously for the finite union. The global-gauge proof gives
the same limits for all physical tests and native poles, including the
prepared places $2,3,17$. Enlarging finite divisibility and the real
slope makes them sufficiently close simultaneously. All native-pole
witnesses are then chosen as before. Consequently replacing independent
first arguments by this progression preserves simultaneous completeness.
It imposes no new condition on an earlier root or reference choice.

\section{The general chart}

Let $\rho\ne0$ be the existing common contact gauge, and retain its
existing registers
\[
             \epsilon=\rho^{-1},\qquad \epsilon_2=\epsilon^2.
\]
Their literal equations are
$\rho\epsilon-1=0$ and $\epsilon_2-\epsilon^2=0$.
For the first contact of block $i$, the three stored coordinates are
\begin{equation}\label{general-scales}
 t_i=\lambda_i/\rho,\qquad
 u_i=(\lambda_i-1)^2,\qquad f_i=(\lambda_i-1)^4.
\end{equation}
They have literal names
$\mathtt{int\_block\_2}$,
$\mathtt{int\_block\_60}$, and
$\mathtt{int\_block\_61}$.
In particular their existing two defining residuals are
\begin{equation}\label{general-rows}
 Q_i=t_i^2-2\epsilon t_i+\epsilon_2-\epsilon_2u_i,
 \qquad P_i=u_i^2-f_i.
\end{equation}
The guards show that these equations have exactly the interpretation
\eqref{general-scales}; they include $\lambda_i=1$.

Retain the three coordinates $t=t_0,u=u_0,f=f_0$ and add seven common
coordinates $d,v,w,a_1,a_2,a_3,a_4$. Impose six new residuals
\begin{align}
 R_v&=\epsilon_2v-d(t-\epsilon),&
 R_w&=\epsilon_2w-d^2,\label{general-mixed}\\
 R_1&=a_1-4uv,&R_2&=a_2-2uw-4v^2,\label{general-fourth}\\
 R_3&=a_3-4vw,&R_4&=a_4-w^2.\notag
\end{align}
Replace the three old coordinates of every $i>0$, in every remaining
residual of the complete system, by the affine expressions
\begin{align}
 t_i&=t+id,\label{affine}\\
 u_i&=u+2iv+i^2w,\notag\\
 f_i&=f+ia_1+i^2a_2+i^3a_3+i^4a_4.\notag
\end{align}
Delete its two rows \eqref{general-rows}, while retaining $Q_0,P_0$.
All new rows and all substituted old rows have degree at most two.

\begin{proposition}[Exact reconstruction]\label{reconstruction}
Every rational solution of the new system reconstructs a solution of
the preceding system by \eqref{affine}. Every preceding solution whose
physical first arguments form an arithmetic progression extends to a
new solution. No guard on the spacing $d$ is required.
\end{proposition}
\begin{proof}
After the substitutions, the omitted rows have the polynomial identities
\begin{equation}\label{identities}
 Q_i=Q_0-2iR_v-i^2R_w,
 \qquad
 P_i=P_0-\sum_{j=1}^4i^jR_j.
\end{equation}
Thus all old residuals hold. These identities are valid as polynomial
identities, independently of any division or nonvanishing assumption.
They also cover $d=0$ and repeated input values.

Conversely, suppose $\lambda_i=\lambda_0+iD$. Set
\[
 d=D/\rho,\qquad v=D(\lambda_0-1),\qquad w=D^2.
\]
Set $a_1=4uv$, $a_2=2uw+4v^2$, $a_3=4vw$, and $a_4=w^2$.
The old guarded equations give \eqref{general-scales}, so direct
expansion gives \eqref{affine} and all six new equations. This proves
the converse. The divisions here are only by the already guarded
$\rho$.
\end{proof}

\section{The specified chart}

The specified interface has different contact weights and reflection
center. Its stored coordinates instead satisfy
\begin{equation}\label{specified-scales}
 t_i=\lambda_i,\quad
 u_i=\rho(22\lambda_i-3)^2,\quad
 f_i=\rho(22\lambda_i-3)^4.
\end{equation}
They have the same local coordinate numbers $2,60,61$. The existing
defining residuals are
\[
 Q_i=(22t_i-3)^2-\epsilon u_i,
 \qquad P_i=u_i^2-\rho f_i.
\]
Keep the same affine substitutions \eqref{affine}, but use
\begin{align}
 R_v&=\epsilon v-22d(22t-3),&
 R_w&=\epsilon w-484d^2,\label{specified-mixed}\\
 R_1&=\rho a_1-4uv,&
 R_2&=\rho a_2-2uw-4v^2,\label{specified-fourth}\\
 R_3&=\rho a_3-4vw,&
 R_4&=\rho a_4-w^2.\notag
\end{align}
Again every residual has degree at most two. The two identities
\eqref{identities} hold literally with these definitions, proving
rational reconstruction of every omitted row.

For the converse use $d=D$, and write $A=22\lambda_0-3$. Set
\[
 v=22\rho DA,\qquad w=484\rho D^2,
\]
and define $a_1,\ldots,a_4$ by \eqref{specified-fourth}. The already
retained guard permits division by $\rho$. Expanding the square and
fourth power of $A+22iD$ proves \eqref{affine}. This also includes
the centered value $A=0$ and does not require $D\ne0$ for soundness.

\section{Complete quartic systems and verification}

For $k$ blocks the transformation deletes $3(k-1)$ coordinates and
$2(k-1)$ residuals, and adds seven coordinates and six residuals.
Its respective savings are $3k-10$ and $2k-8$. Here $k=12$.

\begin{theorem}\label{main}
Under the same arithmetic premises and admissible fixed defining
coefficients as the corrected predecessors, the new complete systems
have the same outer-input projection. Their counts are
\begin{center}
\begin{tabular}{lrrrr}
\toprule
Interface & Old variables & New variables & Old rows & New rows\\
\midrule
General &1074&1048&924&908\\
Specified &1254&1228&1080&1064\\
\bottomrule
\end{tabular}
\end{center}
The sum of squares of the displayed quadratic residuals is an integer
polynomial of total degree exactly four in the counted rational
variables. The integer instance is encoded in its coefficients and is
not counted as a rational witness coordinate.
\end{theorem}
\begin{proof}
For soundness, Proposition~\ref{reconstruction} and its specified
analogue reconstruct the complete preceding rational system. Its
established outer reconstruction gives a witness for the corrected
positive-integer source. All old denominator and nonzero guards remain.

For completeness, start with a positive-integer source witness. Carry
out the existing simultaneous point preparation, then choose $D$ by
Lemma~\ref{spacing}. The argument in the first section supplies the old
complete rational certificates with $\lambda_i=\lambda_0+iD$.
Proposition~\ref{reconstruction} and its specified analogue fill the
new seven registers. The ceiling-power and corrected Pell blocks are
unchanged. Thus the outer-input projection is preserved, although an
arbitrary old certificate need not itself have equally spaced contacts.

The literal census gives the displayed counts. Every row has degree at
most two and coefficients in the integer polynomial ring of the fixed
coefficient parameters. Over $\Q$, a sum of squares vanishes exactly
when every summand vanishes. At least one retained row has nonzero
quadratic homogeneous part, so its sum of squares has degree exactly
four.
\end{proof}

The positive coefficient parameters are $x,z,u,y$, with $x>2z$ on
the retained cofinite instance set; the general table also lists
$C_{\rm parameter}$. After the usual fixed choice of $z,u,y$ from
Jones's coding theorem, the distinguished integer instance is $x$.
The corrected source theorem \cite{corrected} derives its universality
from Jones's classical Theorem~2 \cite[p.~860]{Jones80}, together with
the separately proved binomial and signed-Pell certificates. The old
unresolved expanded-source premise $\mathcal U_{\rm em}$ is not used.
The classical coding theorem remains an external theorem, and the
arithmetic premises of the rational integer interfaces remain explicit
hypotheses. The counts above are therefore qualified by those retained
premises; they do not assert an unconditional undecidability result or
a global minimum.

The complete coefficient files are
\begin{center}\small
\nolinkurl{undecidability/general_contact_v94/affine_contact_m163_quadrics.json}\\
\nolinkurl{undecidability/general_contact_v94/affine_contact_m179_quadrics.json}.
\end{center}
Their ordered coordinate dictionaries, integer coefficient strings,
monomial lists, and ordered residual arrays specify every coefficient
of the final polynomial by the rule $P=\sum_r r^2$. No unlisted
helper, denominator, or arithmetic predicate is implicit in this rule.
The adjacent certificates record the predecessor and output SHA-256
hashes, the exact coordinate and row changes, and the qualification
ledger.

The compiler expands affine substitutions by exact rational coefficient
dictionaries. The independent checker does not import the compiler: it
reconstructs the substitutions symbolically, proves all twenty-two
omitted-row identities, compares every retained row, pins both gauge
guards, and checks integer coefficients, degree, support, and counts.
Run
\begin{center}\small
\texttt{python undecidability/general\_contact\_v94/build\_affine\_contacts.py}\\
\texttt{python undecidability/general\_contact\_v94/verify\_affine\_contacts.py}.
\end{center}
These finite verifications certify the complete algebraic transformation;
the simultaneous arithmetic completeness assertion is the proof above.

\section{Composition with automatic source positivity}

The signed-binomial theorem \cite{signed} applies to the corrected
source with $M=n^2\ge16$ and $2\le b\le2M$. It proves that the
integer inputs $w,s$ are positive from the retained Pell equations,
the strict source inequality, and the single guard $s\ne0$.
Consequently its complete-table transformation deletes the two rows
\texttt{jones1980.positive.w} and
\texttt{jones1980.positive.s}, together with their six private
coordinates, and adds one rational coordinate $s_{\rm inv}$ and
the quadratic residual
\[
                            s s_{\rm inv}-1.
\]
This saves five coordinates and one row. The quantified proof covers
all four signs of $w,s$ and establishes $c\ne0$ before the reciprocal
integer predicate is used; its nonzero and sign deductions do not
depend on the contact charts modified here.

\begin{proposition}[Composed complete systems]
Under the same hypotheses as Theorem~\ref{main}, the complete signed
affine-contact systems have $1043$ variables and $907$ quadratic rows
in the general interface, and $1223$ variables and $1063$ rows in the
specified interface. They have the same outer-input projection as
the corrected ceiling-power predecessors. Their respective savings
from those predecessors are $31$ coordinates and $17$ rows.
\end{proposition}
\begin{proof}
Every coordinate replaced by the contact transformation belongs to an
integer-certificate chart, with a literal name beginning
\texttt{int\_}. Every omitted contact row begins
\texttt{integer.}. By contrast, the six removed positivity coordinates
occur only in their two named source rows. None is an input coordinate
of an integer certificate, and none is replaced in \eqref{affine}.
Thus the two transformations have disjoint changed supports.

From a zero of the composed system, first apply the signed-binomial
theorem. Its hypotheses are available: each rational integer
certificate reconstructs its source input by the unchanged soundness
argument, using the exact contact reconstruction proved above.
The theorem makes $w,s$ positive integers. Their existing
three-triangular-number representations restore the six private
positivity coordinates and the two source rows without altering any
retained coordinate. The contact reconstruction then gives a complete
zero of the corrected ceiling-power predecessor.

Conversely, start with a corrected positive-integer source witness.
Theorem~\ref{main} supplies all the affine-contact certificates.
The positive value of $s$ permits $s_{\rm inv}=1/s$, so dropping
the six now redundant coordinates yields a zero of the composed
system. No new arithmetic premise or coding theorem is used.
The literal savings add: $(26,16)+(5,1)=(31,17)$.
All residuals remain quadratic, with at least one nonzero quadratic
part, so their sum of squares has degree exactly four.
\end{proof}

For comparison, the ordinary rational $abc$ branch applies the same
signed-source theorem after its fixed contact offsets and has $781$
variables and $684$ rows. No general or specified contact progression
is used in that branch.

The two new complete coefficient files are
\begin{center}\small
\nolinkurl{undecidability/general_contact_v94/signed_affine_contact_m163_quadrics.json}\\
\nolinkurl{undecidability/general_contact_v94/signed_affine_contact_m179_quadrics.json}.
\end{center}
Their adjacent certificates record both predecessor hashes and both
stages of the savings. Reproduce and audit them with
\begin{center}\small
\texttt{python undecidability/general\_contact\_v94/build\_signed\_affine\_contacts.py}\\
\texttt{python undecidability/general\_contact\_v94/verify\_signed\_affine\_contacts.py}.
\end{center}
The composite verifier imports neither builder. It compares every
retained residual, independently checks the two deleted positivity
rows and their private supports, pins the new inverse equation,
audits the disjoint changed supports, and verifies coefficient domains,
degree, full coordinate support, counts, qualifications, and hashes.
The quantified source positivity theorem and simultaneous arithmetic
completeness remain the separate mathematical proofs cited above.

\begin{thebibliography}{9}
\bibitem{global}
ZQ Project, \emph{One contact gauge in a combined quadratic system},
\nolinkurl{undecidability/papers/global_gauge_quartic_checkpoint.tex},
especially the ordered simultaneous completeness argument.
\bibitem{common}
ZQ Project, \emph{Common point coordinates in the quartic transfer},
\nolinkurl{undecidability/papers/common_points_quartic_checkpoint.tex},
the common input numerator theorem and all-sign common-index CRT proof.
\bibitem{ceiling}
ZQ Project, \emph{A quadratic ceiling-binary circuit for a fixed power},
\nolinkurl{undecidability/papers/ceiling_power_quartic_checkpoint.tex}.
\bibitem{corrected}
ZQ Project, \emph{A corrected quartic source from Jones's binomial theorem},
\nolinkurl{undecidability/source_repair_v93/corrected_jones_source.tex}.
\bibitem{signed}
ZQ Project, \emph{Signed binomial witnesses force their own positivity
and save five quadratic coordinates},
\nolinkurl{undecidability/source_repair_v93/signed_binomial_positivity.tex}.
\bibitem{Jones80}
J. P. Jones, \emph{Undecidable Diophantine equations},
Bull. Amer. Math. Soc. (N.S.) \textbf{3} (1980), 859--862,
Theorem~2, p.~860, doi:10.1090/S0273-0979-1980-14832-6.
\end{thebibliography}
\end{document}
